\pdfminorversion=7
\documentclass[aspectratio=169,11pt]{beamer}
\usepackage[T1]{fontenc}
\usepackage[utf8]{inputenc}
\usepackage{lmodern,microtype}
\usepackage[english]{babel}
\usepackage{csquotes,amsmath,amssymb,booktabs,array,listings}
\usepackage{appendixnumberbeamer}
\usetheme{Madrid}
\usefonttheme{serif}
\useoutertheme[subsection=false]{miniframes}
\setbeamertemplate{itemize items}[circle]
\setbeamertemplate{enumerate items}[default]
\setbeamertemplate{blocks}[rounded][shadow=false]
\setbeamercolor{alerted text}{fg=red!70!black}
\setbeamerfont{section in head/foot}{size=\scriptsize}
\setbeamertemplate{mini frame in other section}{\usebeamertemplate{mini frame in current subsection}}
\setbeamertemplate{mini frame in other subsection}{\usebeamertemplate{mini frame in current subsection}}
\setbeamertemplate{footline}{%
 \hfill\usebeamerfont{page number in head/foot}\usebeamercolor[fg]{structure}%
 \insertframenumber{} / \inserttotalframenumber\hspace*{2ex}\vskip3pt}
\graphicspath{{assets/}{../exercises/hansen_results/}}
\lstdefinestyle{stata}{basicstyle=\ttfamily\footnotesize,breaklines=true,
 columns=fullflexible,keepspaces=true,showstringspaces=false,
 aboveskip=0.5em,belowskip=0.5em,commentstyle=\color{black!55}}
\lstset{style=stata}
\usepackage[backend=biber,style=authoryear,maxcitenames=2,mincitenames=1,
 maxbibnames=99,giveninits=true,uniquename=false,uniquelist=false,
 doi=false,url=false,isbn=false,eprint=false]{biblatex}
\addbibresource{references.bib}
\DeclareDelimFormat{nameyeardelim}{\addcomma\space}
\DeclareDelimFormat[textcite]{nameyeardelim}{\addspace}
\AtEveryBibitem{\clearlist{language}\clearfield{month}}
\setbeamertemplate{bibliography item}{}
\renewcommand*{\bibfont}{\scriptsize}
\newcommand{\slidecite}[1]{{\footnotesize\color{black!60}\parencite{#1}}}
\newcommand{\slidetextcite}[1]{{\footnotesize\color{black!60}\textcite{#1}}}
\newcommand{\source}[1]{\par\vspace{0.3em}{\scriptsize\color{black!60}#1}}
\newcommand{\E}{\mathbb{E}}
\newcommand{\ind}{\mathbf{1}}
\newcommand{\graph}[2][0.73]{\centering\includegraphics[width=0.96\textwidth,height=#1\textheight,keepaspectratio]{#2}\par}
\title[Regression Discontinuity]{Regression Discontinuity Designs}
\subtitle{Sharp and fuzzy designs, with a Stata exercise}
\author[Luis P. de la Horra]{Luis P.\ de la Horra}
\institute{Universidad de Valladolid}
\date{RDD seminar}

\begin{document}
\begin{frame}[plain]
 \titlepage
 \vspace{-0.5em}\centering
 {\footnotesize Based on Scott Cunningham's \emph{Causal Inference: The Remix} (second edition)}
\end{frame}

\section{Intuition}\subsection{Overview}
\begin{frame}{Seminar outline}
 \begin{enumerate}
  \setlength{\itemsep}{0.6em}
  \item Recognize a cutoff-based research design
  \item Understand continuity and the sharp RD estimand
  \item Estimate the discontinuity and assess credibility
  \item Connect fuzzy RD to instrumental variables
  \item Apply the method to DUI sanctions and recidivism in Stata
 \end{enumerate}
 \vspace{0.6em}
 \textbf{Presentation:} 50--55 minutes \qquad \textbf{Exercise:} 20 minutes
 \source{Reading: Cunningham, Chapters 6 and 7 (fuzzy RDD subsection) \slidecite{cunningham}}
\end{frame}

\begin{frame}{A rule can create a research design}
 \begin{itemize}\setlength{\itemsep}{0.7em}
  \item A score $X_i$ determines eligibility at a known threshold $c$.
  \item Compare the outcome \emph{limits} immediately above and below $c$.
  \item The score can predict outcomes strongly; identification rests on what happens \emph{at the cutoff}.
 \end{itemize}
 \vspace{0.7em}
 \begin{tabular}{@{}ll@{}}\toprule
  Assignment rule & Application in the manual \\\midrule
  Test-score threshold & Merit awards \slidecite{thistlethwaite1960}\\
  Enrollment thresholds & Class size \slidecite{angristlavy1999}\\
  School boundaries & Housing prices \slidecite{black1999}\\\bottomrule
 \end{tabular}\par
 \source{Cunningham, \S6.1; overview: \slidetextcite{leelemieux2010}}
\end{frame}

\begin{frame}{Legal access to alcohol at age 21}
 \graph[0.68]{alcohol_drinking.png}
 {\small Dots: age-bin averages. Curves: separate fits. The jump occurs at age 21.}
 \source{Cunningham, Figure 6.2; original study: \slidetextcite{carpenter2009}}
\end{frame}

\begin{frame}{Mortality at the same threshold}
 \graph[0.67]{alcohol_mortality.png}
 {\small Mortality also jumps. The sharp treatment is \textbf{legal access}, not alcohol consumed.}
 \source{Cunningham, Figure 6.3; \slidecite{carpenter2009}}
\end{frame}

\begin{frame}{Sharp and fuzzy assignment}
 \begin{columns}[T]
  \begin{column}{0.59\textwidth}\graph[0.65]{sharp_fuzzy.pdf}\end{column}
  \begin{column}{0.38\textwidth}
   \textbf{Sharp RD}\\[0.4em]
   $D_i=\ind\{X_i\geq c\}$\\[1em]
   Treatment changes from 0 to 1.\\[1.2em]
   \textbf{Fuzzy RD}\\[0.4em]
   $\Pr(D_i=1\mid X_i=x)$ jumps, with imperfect take-up.
  \end{column}
 \end{columns}
 \source{Cunningham, Figure 6.4; adapted there from \slidetextcite{vanderklaauw2002}}
\end{frame}

\section{Sharp RD}\subsection{Identification}
\begin{frame}{Potential outcomes and observed outcomes}
 \begin{itemize}\setlength{\itemsep}{0.6em}
  \item Define treatment first: $D_i=1$ means receiving the intervention.
  \item $Y_i(1)$ and $Y_i(0)$ exist for every unit, whatever its actual treatment.
  \item We observe only $Y_i=D_iY_i(1)+(1-D_i)Y_i(0)$.
 \end{itemize}
 \vspace{0.6em}
 \begin{tabular}{@{}lll@{}}\toprule
  Sharp assignment & Observed & Missing counterfactual \\\midrule
  $X_i<c$ & $Y_i(0)$ & $Y_i(1)$ \\
  $X_i\geq c$ & $Y_i(1)$ & $Y_i(0)$ \\\bottomrule
 \end{tabular}\par
 \vspace{0.6em}
 No treated and untreated units share the same score in a sharp design.
 \source{Cunningham, \S6.2; supplied smoothness handout, pp. 1--3 \slidecite{cunningham}}
\end{frame}

\begin{frame}{Two potential-outcome curves}
 \begin{columns}[T]
  \begin{column}{0.68\textwidth}\graph[0.68]{smoothness.png}\end{column}
  \begin{column}{0.29\textwidth}\small
   $m_d(x)=\E[Y_i(d)\mid X_i=x]$\\[0.8em]
   Hold treatment fixed as you move horizontally.\\[0.8em]
   Each curve describes group averages, not one person's trajectory.\\[0.8em]
   Both schedules are visible here because the data are simulated.
  \end{column}
 \end{columns}
 \source{Cunningham, Figure 6.7; supplied smoothness handout \slidecite{hahn2001}}
\end{frame}

\begin{frame}{The continuity assumption}
 For each fixed treatment status $d\in\{0,1\}$:
 \[
  \lim_{x\to c^-}\E[Y_i(d)\mid X_i=x]
  =\lim_{x\to c^+}\E[Y_i(d)\mid X_i=x]=m_d(c).
 \]
 \begin{block}{Read horizontally}
  Applicants on opposite sides approach the same mean outcome under the \emph{same hypothetical treatment status}.
 \end{block}
 \begin{itemize}\setlength{\itemsep}{0.45em}
  \item Each potential-outcome mean is continuous at $c$.
  \item Curves may slope, bend, and have different slopes.
  \item Continuity neither implies a zero treatment effect nor equal individual outcomes.
 \end{itemize}
 \source{Cunningham, \S6.2; supplied handout, p. 2 \slidecite{hahn2001}}
\end{frame}

\begin{frame}{Why the observed outcome can jump}
 \graph[0.60]{smoothness_observed_counterfactual.pdf}
 \begin{itemize}\setlength{\itemsep}{0.3em}
  \item The observed mean follows $m_0(x)$ below $c$ and $m_1(x)$ above $c$.
  \item The vertical gap at $c$ is $m_1(c)-m_0(c)$.
 \end{itemize}
 \source{Supplied smoothness handout, p. 3; schematic values are illustrative \slidecite{hahn2001}}
\end{frame}

\begin{frame}{Identification at the cutoff}
 \begin{align*}
  \tau_{\mathrm{SRD}}(c)
   &=\E[Y_i(1)-Y_i(0)\mid X_i=c]\\[0.6em]
   &=\underbrace{\lim_{x\to c^+}\E[Y_i\mid X_i=x]}_{m_1(c)}
     -\underbrace{\lim_{x\to c^-}\E[Y_i\mid X_i=x]}_{m_0(c)}.
 \end{align*}
 \begin{itemize}\setlength{\itemsep}{0.6em}
  \item Assignment tells us which outcome we observe on each side.
  \item Continuity connects those observed limits to the missing means at $c$.
  \item This identifies the \textbf{ATE at the cutoff}.
 \end{itemize}
 \source{Cunningham, \S6.2; \slidecite{hahn2001,imbenslemieux2008}}
\end{frame}

\begin{frame}{A local effect need not generalize}
 \graph[0.65]{potential_heterogeneous.pdf}
 {\small Continuous curves can imply different effects at different scores. RD identifies the gap at $c$.}
 \source{Cunningham, Figure 6.6 \slidecite{cunningham,leelemieux2010}}
\end{frame}

\begin{frame}{A university example for the counterfactuals}
 \textbf{Teaching simplification:} everyone above the admission cutoff attends the flagship.
 \[
  X_i=\mathrm{SAT}_i-\mathrm{required\ SAT}(\mathrm{GPA}_i,\mathrm{cohort}_i),\qquad c=0.
 \]
 \begin{itemize}\setlength{\itemsep}{0.65em}
  \item $Y_i(1)$: earnings if applicant $i$ attends the flagship.
  \item $Y_i(0)$: earnings if the same applicant does not attend it; another university is possible.
  \item Imagine \emph{nobody} attended: mean earnings should pass continuously through zero.
  \item Imagine \emph{everybody} attended: that mean should also pass continuously through zero.
 \end{itemize}
 \source{Supplied handout, pp. 4--5. The actual study is fuzzy RD \slidecite{hoekstra2009}}
\end{frame}

\begin{frame}{Continuity with a positive treatment effect}
 \centering
 \renewcommand{\arraystretch}{1.4}
 \begin{tabular}{@{}lrr@{}}\toprule
  Adjusted SAT & Without flagship & With flagship\\\midrule
  $-10$ & \$39,650 (observed) & \$49,500 (missing)\\
  $0$ & \$40,000 (missing) & \$50,000 (observed)\\
  $+10$ & \$40,350 (missing) & \$50,500 (observed)\\\bottomrule
 \end{tabular}\par
 \vspace{1em}\raggedright
 Both schedules are continuous. The cutoff effect is \textbf{\$10,000}.
 \vspace{0.7em}
 The observed difference between $-10$ and $+10$ is \$10,850: it also reflects differences associated with score.
 \source{Invented dollar amounts from the supplied handout. Hoekstra analyzes log earnings; these are not paper estimates \slidecite{hoekstra2009}}
\end{frame}

\section{Estimation}\subsection{Methods}
\begin{frame}[fragile]{OLS with separate slopes}
 Center the running variable: $x_i=X_i-c$.
 \[
  Y_i=\alpha+\tau D_i+\beta x_i+\gamma D_ix_i+\varepsilon_i.
 \]
 \begin{itemize}\setlength{\itemsep}{0.4em}
  \item The left intercept is $\alpha$; the right intercept is $\alpha+\tau$.
  \item The slopes are $\beta$ and $\beta+\gamma$.
  \item $\hat\tau$ is the fitted gap \emph{at the cutoff}.
 \end{itemize}
 \begin{lstlisting}
generate double x_c = score - cutoff
generate byte treated = (score >= cutoff) if !missing(score)
regress outcome i.treated##c.x_c if abs(x_c) <= h, vce(robust)
 \end{lstlisting}
 \source{Generic Stata syntax; $h$ is a chosen window. Cunningham, \S6.3 \slidecite{cunningham}}
\end{frame}

\begin{frame}{Functional form can create a false jump}
 \begin{columns}[T]
  \begin{column}{0.60\textwidth}\graph[0.62]{misspecification.png}\end{column}
  \begin{column}{0.37\textwidth}\small
   The manual simulates \textbf{zero treatment effect} with a nonlinear outcome-score relationship.\\[0.8em]
   Broad-window straight lines give misleading cutoff predictions.\\[0.8em]
   Adding high global polynomial orders is unstable; use local low-order fits.
  \end{column}
 \end{columns}
 \source{Cunningham, Figure 6.9 and Table 6.1 \slidecite{gelman2019}}
\end{frame}

\begin{frame}{Bin averages and boundary bias}
 \graph[0.58]{boundary_bias.pdf}
 {\small A bin average summarizes outcomes away from the cutoff. With trends, its difference need not equal the boundary gap.}
 \vspace{0.4em}
 \textbf{Binning is for display; estimation uses individual observations.}
 \source{Cunningham, Figure 6.10; binscatter discussion: \slidetextcite{binscatter2024}}
\end{frame}

\begin{frame}{Local linear regression}
 Fit a weighted regression separately on each side:
 \[
  (\hat a_s,\hat b_s)=\arg\min_{a,b}
   \sum_{i\in s,\,|x_i|\leq h_s}K(x_i/h_s)(Y_i-a-bx_i)^2.
 \]
 \[
  \hat\tau=\hat a_+-\hat a_-,\qquad
  K(u)=(1-|u|)\ind\{|u|\leq1\}.
 \]
 \begin{itemize}\setlength{\itemsep}{0.6em}
  \item \textbf{Local:} use a neighborhood around $c$.
  \item \textbf{Linear:} allow different slopes on either side.
  \item \textbf{Triangular weights:} give nearby observations more influence.
 \end{itemize}
 \source{Cunningham, \S6.3 \slidecite{hahn2001,calonico2014,cattaneo2020}}
\end{frame}

\begin{frame}{Bandwidth, kernel, and polynomial order}
 \begin{itemize}\setlength{\itemsep}{0.65em}
  \item \textbf{Bandwidth $h$:} a narrower window usually reduces approximation bias but uses less information.
  \item \textbf{Kernel $K$:} determines observation weights inside the window.
  \item \textbf{Order $p$:} local linear ($p=1$) is a standard starting point.
 \end{itemize}
 \vspace{0.7em}
 \[
  \operatorname{MSE}(\hat\tau)=\operatorname{Bias}(\hat\tau)^2+\operatorname{Var}(\hat\tau).
 \]
 Data-driven bandwidths balance bias and variance; they can differ by side.
 \vspace{0.5em}
 Report the chosen windows and examine sensitivity.
 \source{Cunningham, \S6.3 \slidecite{imbenskalyanaraman2012,calonico2014}}
\end{frame}

\begin{frame}{Robust bias-corrected inference}
 At an MSE-optimal bandwidth, approximation bias can matter for confidence intervals.
 \par
 \vspace{0.9em}
 \centering
 \renewcommand{\arraystretch}{1.5}
 \begin{tabular}{@{}lll@{}}\toprule
  Interval & Center & Standard error\\\midrule
  Conventional & $\hat\tau$ & Conventional\\
  Bias-corrected & $\hat\tau-\hat B$ & Conventional\\
  Robust bias-corrected & $\hat\tau-\hat B$ & Includes bias estimation\\\bottomrule
 \end{tabular}\par
 \vspace{0.9em}\raggedright
 \[
  \mathrm{CI}_{\mathrm{RBC}}=(\hat\tau-\hat B)\pm1.96\,\mathrm{SE}_{\mathrm{RBC}}.
 \]
 In \texttt{rdrobust}, pair \texttt{e(tau\_bc)} with \texttt{e(se\_tau\_rb)}.
 \source{Cunningham, \S6.3 \slidecite{calonico2014,calonico2017}}
\end{frame}

\section{Validity}\subsection{Evidence}
\begin{frame}{What can invalidate the comparison?}
 \begin{itemize}\setlength{\itemsep}{0.7em}
  \item Precise sorting: people manipulate their scores around $c$.
  \item Another intervention changes at exactly the same cutoff.
  \item Measurement, sample selection, or reporting rules change there.
  \item Nonrandom heaping makes neighboring scores different types of units.
 \end{itemize}
 \vspace{0.7em}
 \begin{block}{A failure of continuity}
  If untreated potential outcomes jump at $c$, the observed discontinuity can combine that jump with the treatment effect.
 \end{block}
 \source{Cunningham, \S6.4 \slidecite{imbenslemieux2008,leelemieux2010}}
\end{frame}

\begin{frame}{Density continuity and manipulation}
 \graph[0.59]{density.jpg}
 \begin{itemize}\setlength{\itemsep}{0.3em}
  \item Look for a change in the distribution of scores at $c$.
  \item A density test is evidence about sorting; it cannot verify missing potential outcomes.
 \end{itemize}
 \source{Cunningham, Figure 6.11 \slidecite{mccrary2008,cjm2020}}
\end{frame}

\begin{frame}{Predetermined covariates and placebo cutoffs}
 \begin{columns}[T]
  \begin{column}{0.62\textwidth}\graph[0.66]{election_balance.jpg}\end{column}
  \begin{column}{0.35\textwidth}\small
   In close elections, check whether district characteristics jump at 50\% Democratic vote share.\\[0.8em]
   Also estimate jumps at placebo cutoffs within each side.\\[0.8em]
   Null findings support credibility; they do not establish continuity.
  \end{column}
 \end{columns}
 \source{Cunningham, Figure 6.22 (covariate checks) \slidecite{leemoretti2004,imbenslemieux2008}}
\end{frame}

\begin{frame}{Birth-weight heaping}
 \begin{columns}[T]
  \begin{column}{0.61\textwidth}\graph[0.64]{birthweight_heaping.jpg}\end{column}
  \begin{column}{0.36\textwidth}\small
   Very-low-birth-weight classification changes at 1,500 grams.\\[0.7em]
   Reported weights heap at round values.\\[0.7em]
   If heaped babies differ systematically, the local comparison can be biased.\\[0.7em]
   Excluding the cutoff changes the sample and needs justification.
  \end{column}
 \end{columns}
 \source{Cunningham, Figure 6.13 \slidecite{almond2010,barreca2011,barreca2016}}
\end{frame}

\begin{frame}{Medicare eligibility at age 65}
 \begin{columns}[T]
  \begin{column}{0.60\textwidth}\graph[0.62]{medicare_employment.jpg}\end{column}
  \begin{column}{0.37\textwidth}\small
   Age 65 changes eligibility sharply; actual insurance coverage has imperfect take-up.\\[0.8em]
   Retirement is a potential competing mechanism.\\[0.8em]
   The manual uses employment profiles to assess whether work changes abruptly at 65.
  \end{column}
 \end{columns}
 \source{Cunningham, \S6.5 and Figure 6.14 \slidecite{card2008}}
\end{frame}

\begin{frame}{Healthcare utilization at age 65}
 \graph[0.65]{medicare_hospital.jpg}
 {\small This is a reduced-form effect of crossing the eligibility threshold. An effect of actual coverage requires a first stage and additional assumptions.}
 \source{Cunningham, Figure 6.15; utilization and subsequent mortality study \slidecite{card2008,card2009}}
\end{frame}

\section{Fuzzy RD}\subsection{Imperfect take-up}
\begin{frame}{Hoekstra's actual design: admission and attendance}
 \begin{columns}[T]
  \begin{column}{0.59\textwidth}\graph[0.65]{hoekstra_firststage.jpg}\end{column}
  \begin{column}{0.38\textwidth}\small
   $Z_i=\ind\{X_i\geq0\}$: crossing the admission threshold.\\[0.7em]
   $D_i$: attending the flagship.\\[0.7em]
   Admission causes a jump in enrollment, rather than universal attendance.\\[0.7em]
   Unpublished, changing cutoffs make precise targeting less plausible.
  \end{column}
 \end{columns}
 \source{Cunningham, Figure 6.20; supplied handout, pp. 4--5 \slidecite{hoekstra2009}}
\end{frame}

\begin{frame}{The fuzzy RD ratio}
 \[
  \tau_{\mathrm{FRD}}(c)=
  \frac{\lim_{x\to c^+}\E[Y_i\mid X_i=x]-\lim_{x\to c^-}\E[Y_i\mid X_i=x]}
       {\lim_{x\to c^+}\E[D_i\mid X_i=x]-\lim_{x\to c^-}\E[D_i\mid X_i=x]}.
 \]
 \vspace{0.6em}
 \begin{itemize}\setlength{\itemsep}{0.7em}
  \item \textbf{Numerator:} reduced-form jump in the outcome.
  \item \textbf{Denominator:} first-stage jump in treatment take-up.
  \item \textbf{Illustration:} an outcome jump of 2 and a take-up jump of 0.4 imply a ratio of 5.
 \end{itemize}
 \source{Illustrative numbers. Cunningham, \S7.7, fuzzy RDD subsection \slidecite{hahn2001,imbenslemieux2008}}
\end{frame}

\begin{frame}{Additional assumptions and the estimand}
 \begin{itemize}\setlength{\itemsep}{0.65em}
  \item \textbf{Continuity:} the relevant potential outcomes and compliance types vary smoothly at $c$.
  \item \textbf{Relevance:} crossing the cutoff changes the probability of treatment.
  \item \textbf{Exclusion:} crossing the cutoff affects $Y$ through $D$.
  \item \textbf{Monotonicity:} crossing to the eligible side does not reduce anyone's take-up.
  \item \textbf{Well-defined treatment and no interference.}
 \end{itemize}
 \vspace{0.6em}
 The ratio identifies the \textbf{LATE for compliers at the cutoff}.
 \source{Cunningham, fuzzy RDD subsection \slidecite{hahn2001,imbensangrist1994,angrist1996}}
\end{frame}

\begin{frame}[fragile]{Fuzzy RD as local 2SLS}
 With $x_i=X_i-c$ and $Z_i=\ind\{X_i\geq c\}$:
 \begin{align*}
 D_i&=a+\pi Z_i+b x_i+g Z_ix_i+v_i,\\
 Y_i&=\alpha+\tau D_i+\beta x_i+\gamma Z_ix_i+u_i.
 \end{align*}
 Instrument $D_i$ with $Z_i$; retain the same score controls on each side.
 \begin{lstlisting}
generate double zx = z*x_c
ivregress 2sls earnings x_c zx (attend = z) ///
    if abs(x_c) <= h, vce(robust)
* Alternatively: local polynomial fuzzy RD with RBC inference
rdrobust earnings score, c(0) fuzzy(attend) p(1)
 \end{lstlisting}
 \source{Generic syntax, not a run on Hoekstra data. Cunningham, fuzzy RDD subsection \slidecite{calonico2017}}
\end{frame}

\begin{frame}{Another fuzzy design: Maimonides' rule}
 \textbf{Angrist and Lavy:} classes in Israeli schools are split at enrollment thresholds. \slidecite{angristlavy1999}
 \vspace{0.7em}
 \begin{itemize}\setlength{\itemsep}{0.6em}
  \item Rule-predicted class size falls when enrollment crosses a threshold.
  \item Actual class size responds imperfectly, creating a first stage.
  \item Outcomes are student achievement measures.
  \item The IV logic asks whether the rule changes achievement through class size.
 \end{itemize}
 \vspace{0.7em}
 Class size is a \emph{continuous} treatment: the binary-treatment complier interpretation requires adaptation.
 \source{Cunningham, \S6.1--6.2 and fuzzy RDD subsection \slidecite{cunningham}}
\end{frame}

\begin{frame}{Sharp and fuzzy RD: what is identified?}
 \renewcommand{\arraystretch}{1.5}
 \begin{tabular}{@{}>{\raggedright\arraybackslash}p{0.21\textwidth}>{\raggedright\arraybackslash}p{0.34\textwidth}>{\raggedright\arraybackslash}p{0.37\textwidth}@{}}\toprule
  & Sharp RD & Fuzzy RD\\\midrule
  Assignment & Treatment is determined by crossing $c$ & Treatment probability changes at $c$\\
  Estimator & Outcome discontinuity & Outcome jump / take-up jump\\
  Estimand & ATE at $c$ & Complier LATE at $c$\\
  Main requirements & Continuity, well-defined treatment & Continuity plus IV assumptions\\\bottomrule
 \end{tabular}\par
 \vspace{0.8em}
 For either design, define treatment precisely and inspect the institutional rule.
 \source{Cunningham, Chapters 6--7 \slidecite{hahn2001,leelemieux2010}}
\end{frame}

\section{Hansen lab}\subsection{Stata exercise}
\begin{frame}{Exercise: DUI sanctions and recidivism}
 \textbf{Research question:} does crossing the DUI BAC threshold reduce subsequent offending? \slidecite{hansen2015}
 \vspace{0.5em}
 \begin{itemize}\setlength{\itemsep}{0.55em}
  \item $X_i$: measured blood alcohol content (BAC); cutoff $c=0.08$.
  \item $D_i=\ind\{X_i\geq0.08\}$: exposure to the threshold-based DUI regime.
  \item $Y_i$: recidivism indicator in the teaching extract.
  \item Cunningham's extract has \textbf{214,558 observations}; this is a teaching exercise rather than a full replication of the paper.
 \end{itemize}
 \vspace{0.4em}
 \textbf{20-minute task:} assess credibility, estimate the gap, and interpret its magnitude.
 \source{Cunningham, \S6.6; sanctions and deterrence framework \slidecite{cunningham,becker1968}}
\end{frame}

\begin{frame}[fragile]{Load the data and center BAC}
 Run directly in Stata 18. If necessary, select this do-file in the folder-location dialog:
 \begin{lstlisting}
do "hansen_seminar.do"
 \end{lstlisting}
 The script downloads once and then uses the cached data:
 \begin{lstlisting}
use "`output_dir'/hansen_dwi.dta", clear
rename bac1 bac1_orig
generate double bac1 = bac1_orig - 0.08
generate byte dui = (bac1_orig >= 0.08) if !missing(bac1_orig)
generate double bac1_sq = bac1^2
 \end{lstlisting}
 \begin{itemize}
  \item Use \texttt{bac1\_orig} with \texttt{c(0.08)}; centered BAC has cutoff zero.
  \item The script sets its folder with \texttt{cd}, then saves figures, tables, and a log.
 \end{itemize}
 \source{Seminar script adapted from the supplied \texttt{hansen\_complete.do} \slidecite{hansen2015,cunningham}}
\end{frame}

\begin{frame}[fragile]{Inspect the BAC distribution}
 \begin{lstlisting}
histogram bac1_orig, discrete width(0.001) color(gs10%90) ///
    xline(0.08) title("Figure 6.16: BAC histogram")
 \end{lstlisting}
 \graph[0.48]{figure_6_16.png}
 \source{Generated from the teaching extract; compare Cunningham, Figure 6.16 \slidecite{hansen2015}}
\end{frame}

\begin{frame}[fragile]{The density test and recorded precision}
 \begin{columns}[T]
  \begin{column}{0.60\textwidth}\graph[0.57]{density_diagnostic.png}\end{column}
  \begin{column}{0.37\textwidth}\small
   Robust density-test $p=0.8897$.\\[0.7em]
   It does not reject density continuity.\\[0.6em]
   Narrow-window binomial tests reject at several widths; recorded precision matters.
  \end{column}
 \end{columns}
 \begin{lstlisting}
rddensity bac1_orig, c(0.08) plot
 \end{lstlisting}
 \source{Generated output. Interpret the tests with the histogram \slidecite{mccrary2008,cjm2020}}
\end{frame}

\begin{frame}[fragile]{Covariate balance: code and estimates}
 \begin{lstlisting}
foreach outcome in white male aged acc {
    regress `outcome' i.dui##c.bac1 ///
        if inrange(bac1_orig, 0.03, 0.13), vce(robust)
}
 \end{lstlisting}
 \begin{center}\small
\renewcommand{\arraystretch}{1.15}
\begin{tabular}{@{}lrrrr@{}}\toprule
 & White & Male & Age & Accident\\\midrule
Discontinuity & 0.0057 & 0.0062 & -0.1405 & -0.0034\\
Robust SE & 0.0050 & 0.0057 & 0.1644 & 0.0041\\
Sample mean & 0.8526 & 0.7919 & 34.1701 & 0.1012\\
Untreated cutoff level & 0.8463 & 0.7842 & 33.9199 & 0.0834\\
Observations & 89,967 & 89,967 & 89,967 & 89,967\\
\bottomrule\end{tabular}\end{center}

 \source{Generated Table 6.6; proportions for White/Male/Accident, years for Age \slidecite{hansen2015}}
\end{frame}

\begin{frame}[fragile]{Covariate balance: figures}
 \graph[0.54]{figure_6_17.png}
 \begin{lstlisting}
cmogram male bac1_orig if inrange(bac1_orig, 0.03, 0.13), ///
    scatter cut(0.08) line(0.08) lfit
 \end{lstlisting}
 \source{Repeat for White, Age, and Accident; generated Figure 6.17 \slidecite{hansen2015}}
\end{frame}

\begin{frame}[fragile]{Plot the recidivism discontinuity}
 \begin{lstlisting}
cmogram recidivism bac1_orig if inrange(bac1_orig, 0.03, 0.13), ///
    scatter cut(0.08) line(0.08) lfit
* Repeat with qfit for the quadratic panel.
 \end{lstlisting}
 \graph[0.47]{figure_6_18.png}
 \source{Generated Figure 6.18; read the fitted gap at 0.08 \slidecite{hansen2015}}
\end{frame}

\begin{frame}[fragile]{OLS results in the fixed BAC window}
 \begin{lstlisting}
regress recidivism i.dui##c.bac1 ///
    if inrange(bac1_orig, 0.03, 0.13), vce(robust)
regress recidivism i.dui##c.(bac1 bac1_sq) ///
    if inrange(bac1_orig, 0.03, 0.13), vce(robust)
 \end{lstlisting}
 \begin{center}\small
\renewcommand{\arraystretch}{1.15}
\begin{tabular}{@{}lrr@{}}\toprule
 & Linear & Quadratic\\\midrule
Effect (pp) & -2.362 & -1.432\\
Robust SE (pp) & 0.436 & 0.623\\
Sample recidivism (\%) & 10.694 & 10.694\\
Untreated cutoff level (\%) & 11.707 & 11.134\\
Observations & 89,967 & 89,967\\
\bottomrule\end{tabular}\end{center}

 \source{Generated Table 6.7: sample mean and fitted untreated cutoff level are distinct \slidecite{hansen2015}}
\end{frame}

\begin{frame}[fragile]{Local polynomial RD: code}
 \begin{lstlisting}
rdrobust recidivism bac1_orig, c(0.08) p(1) ///
    kernel(triangular) bwselect(msetwo) ///
    vce(nn 3) masspoints(adjust) all
 \end{lstlisting}
 \begin{itemize}\setlength{\itemsep}{0.65em}
  \item Repeat for $p=0,1,2$: local constant, linear, and quadratic.
  \item Use the full BAC sample for bandwidth selection.
  \item \texttt{msetwo} chooses separate MSE-optimal windows by side.
  \item Pair the bias-corrected effect with the RBC standard error.
 \end{itemize}
 \source{The script stores each model and its bandwidths \slidecite{calonico2014,calonico2017}}
\end{frame}

\begin{frame}{Local polynomial RD: results}
 \begin{center}\small
\renewcommand{\arraystretch}{1.15}
\begin{tabular}{@{}lrrr@{}}\toprule
 & Constant & Linear & Quadratic\\\midrule
Bias-corrected effect (pp) & -1.891 & -1.955 & -2.093\\
RBC SE (pp) & 0.507 & 0.607 & 0.712\\
Polynomial order & 0 & 1 & 2\\
Bandwidth left & 0.0201 & 0.0329 & 0.0378\\
Bandwidth right & 0.0143 & 0.0544 & 0.0938\\
Effective N left & 13,794 & 16,774 & 17,545\\
Effective N right & 16,352 & 77,539 & 135,457\\
Effective N total & 30,146 & 94,313 & 153,002\\
\bottomrule\end{tabular}\end{center}

 \vspace{0.5em}
 \small The quadratic right window reaches BAC 0.174, beyond the aggravated-DUI threshold of 0.15. Check sensitivity to narrower windows.
 \source{Generated Table 6.8: both sides and their sum are reported \slidecite{hansen2015,calonico2017}}
\end{frame}

\begin{frame}{The local polynomial fits}
 \graph[0.51]{table6_8_plots.png}
 \begin{itemize}\setlength{\itemsep}{0.4em}
  \item Each panel uses its own estimated left and right bandwidths.
  \item Curves show conventional fits; subtitles show bias-corrected effects.
  \item Binned dots summarize the data; they do not determine the estimates.
 \end{itemize}
 \source{Generated from the stored Table 6.8 models \slidecite{hansen2015,calonico2017}}
\end{frame}

\begin{frame}{Interpret the local linear result}
 \begin{block}{Local linear estimate}
A reduction of \textbf{1.95 percentage points} in recidivism.\\[0.6em]
Robust bias-corrected 95\% CI: $[-3.14,\ -0.77]$ percentage points.
\end{block}

 \begin{itemize}\setlength{\itemsep}{0.65em}
  \item This is a local effect of crossing the threshold-based DUI regime.
  \item It combines the sanctions and other consequences triggered there.
  \item The data do not separately identify each mechanism behind reduced recidivism.
 \end{itemize}
 \source{Generated $p=1$ result from Cunningham's teaching extract \slidecite{hansen2015}}
\end{frame}

\begin{frame}{Exercise discussion}
 \begin{enumerate}\setlength{\itemsep}{0.7em}
  \item What makes continuity plausible in this institutional setting?
  \item What do the histogram, density test, and balance checks establish?
  \item Why do OLS and local linear estimates differ?
  \item Which people does the estimate describe, and what does it say about mechanism?
 \end{enumerate}
 \vspace{0.7em}
 \begin{block}{Evidence to take away}
  A known rule, a visible outcome gap, credible local comparisons, and inference that accounts for approximation bias.
 \end{block}
 \source{Cunningham, \S6.6--6.7 \slidecite{imbenslemieux2008,leelemieux2010}}
\end{frame}

\appendix
\setbeamertemplate{headline}{}
\section{Supplementary material}

\begin{frame}{Observed data from the manual's simulation}
 \graph[0.70]{observed_jump.png}
 \source{Cunningham, Figure 6.8: observing different potential outcomes on either side \slidecite{cunningham}}
\end{frame}

\begin{frame}{Homogeneous effects as a special case}
 \graph[0.70]{potential_homogeneous.pdf}
 \source{Cunningham, Figure 6.5; broader representativeness requires further assumptions \slidecite{cunningham}}
\end{frame}

\begin{frame}{Discrete scores and standard errors}
 \begin{itemize}\setlength{\itemsep}{0.7em}
  \item Many observations can share a score without providing many distinct score values near $c$.
  \item With fixed support and gaps, a continuous-score limiting argument needs additional structure.
  \item \texttt{masspoints(adjust)} helps bandwidth selection; it does not resolve all identification and inference issues from discreteness.
  \item Clustering only by running-variable values is not a general solution.
 \end{itemize}
 \source{Cunningham, \S6.3; early and later treatments \slidecite{leecard2008,kolesar2018}}
\end{frame}

\begin{frame}[fragile]{Plot the model's actual estimation window}
 \begin{lstlisting}
estimates restore recid_rd_p1
local hl = e(h_l)
local hr = e(h_r)
local xmin = e(c) - `hl'
local xmax = e(c) + `hr'
rdplot recidivism bac1_orig if inrange(bac1_orig, `xmin', `xmax'), ///
    c(0.08) p(1) h(`hl' `hr') kernel(triangular) ///
    nbins(10 10) masspoints(adjust)
 \end{lstlisting}
 \vspace{0.6em}
 This reuses the stored model's bandwidths and the same polynomial order.
 \source{Seminar script, section 7 \slidecite{calonico2017}}
\end{frame}

\begin{frame}{Local constant fit}
 \graph[0.77]{table6_8_plot_p0.png}
 \source{Generated from the Table 6.8 model with $p=0$ \slidecite{hansen2015}}
\end{frame}
\begin{frame}{Local linear fit}
 \graph[0.77]{table6_8_plot_p1.png}
 \source{Generated from the Table 6.8 model with $p=1$ \slidecite{hansen2015}}
\end{frame}
\begin{frame}{Local quadratic fit}
 \graph[0.77]{table6_8_plot_p2.png}
 \source{Generated from the Table 6.8 model with $p=2$ \slidecite{hansen2015}}
\end{frame}

\begin{frame}[fragile]{Store the estimated effects and uncertainty}
 \begin{lstlisting}
forvalues p = 0/2 {
    local col = `p' + 1
    capture program drop rdrobust
    rdrobust recidivism bac1_orig, c(0.08) p(`p') ///
        kernel(triangular) bwselect(msetwo) ///
        vce(nn 3) masspoints(adjust) all
    estimates store recid_rd_p`p'
    matrix table6_8[1, `col'] = e(tau_bc)
    matrix table6_8[2, `col'] = e(se_tau_rb)
}
 \end{lstlisting}
 \source{Excerpt from the runnable seminar script; its matrix also stores bandwidths and effective sample sizes \slidecite{calonico2017}}
\end{frame}

\begin{frame}[fragile]{Export the figures and tables}
 \begin{lstlisting}
graph save fig616 "`output_dir'/figure_6_16.gph", replace
graph export "`output_dir'/figure_6_16.png", ///
    name(fig616) width(2400) replace
* Example for one table:
preserve
clear
svmat double table6_7, names(col)
export delimited using "`output_dir'/table_6_7_values.csv", replace
restore
 \end{lstlisting}
 The full script additionally exports descriptive row labels for all three tables.
 \source{Figures remain editable in Stata's GPH format. The deck's table inputs are generated from CSV output.}
\end{frame}

\begin{frame}{Other deterrence studies in the manual}
 \begin{itemize}\setlength{\itemsep}{0.8em}
  \item An Italian collective pardon changes the sentence faced after reoffending. \slidecite{drago2009}
  \item The adult criminal-justice threshold at age 18 offers another RD setting. \slidecite{leemccrary2017}
  \item These designs distinguish behavioral responses to threatened sanctions from the mechanical effects of incarceration.
 \end{itemize}
 \source{Cunningham, \S6.6; supplementary context for the Hansen exercise \slidecite{becker1968}}
\end{frame}

\begin{frame}[allowframebreaks]{References}
 \printbibliography[heading=none]
\end{frame}
\end{document}
